\documentclass[12pt]{article}

\usepackage{sbc-template}
\usepackage{graphicx,url}
\usepackage[utf8]{inputenc}
\usepackage[T1]{fontenc}
\usepackage[portuguese]{babel}
\usepackage{geometry}
\usepackage{setspace}
\usepackage{titlesec}
     
\sloppy

\title{Uma Análise de NeRFs Híbridas baseadas em CNN e ViT em Conjunto de Dados de Vistas Esparsas}

\author{Felipe G. da Silva\inst{1}}


\address{Instituto de Biociências, Letras e Ciências Exatas -- Universidade Estadual Paulista
  (UNESP)\\
  CEP: 15054-000 -- São José do Rio Preto -- SP -- Brazil
  \email{felipe-gomes.silva@unesp.br}
}

\begin{document} 

\maketitle

\begin{abstract}
 This work compares hybrid NeRFs based on ViT and CNN applied to sparse-view datasets, evaluating the impact of these feature extraction architectures on novel view synthesis performance. Through a quantitative analysis using visual fidelity (PSNR, SSIM, LPIPS) and computational efficiency metrics (convergence time, VRAM footprint, FPS), this study highlights the trade-offs between the global generalization of Transformers and the local processing efficiency of convolutions. Ultimately, it aims to guide the selection of 3D rendering models for practical scenarios often constrained by hardware limitations and data scarcity.
\end{abstract}

\begin{resumo}
  Esse trabalho apresenta uma comparação entre NeRFs híbridas baseadas em ViT e CNN aplicadas em conjuntos de dados com vistas esparsas e avalia o impacto dessas arquiteturas de extração de características no desempenho em síntese de novas vistas. Por meio de uma análise quantitativa baseada em métricas de fidelidade visual (PSNR, SSIM e LPIPS) e de eficiência computacional (tempo de convergência, consumo de VRAM e FPS), o estudo busca evidenciar os trade-offs entre a capacidade de generalização global dos Transformers e a eficiência de processamento local das convoluções,  visando auxiliar na seleção de modelos de renderização 3D em cenários práticos, os quais são frequentemente marcados por restrições de hardware e escassez de dados.
\end{resumo}

\section{Introdução}

Modelos de campos de radiância neural (NeRF) conquistaram, nos anos recentes, um papel muito importante, tanto na indústria, quanto na academia, com resultados hiper-realistas sobre geração de novas vistas \cite{kerbl3Dgaussians, Tancik2023, substance3dviewer}. Entretanto, as mais avançadas técnicas de renderização neural exigem uma grande quantidade de (i) poder computacional; (ii) múltiplas vistas no conjunto de dados. Ambos os desafios supracitados tendem a impedir que aplicações mais simples, como escaneamento \textit{mobile} de objetos, provador virtual, geração de assets 3D e síntese em datasets pequenos, sejam aproveitados. É com tal problemática que projetos como Instant-NGP (\textit{Instant Neural Graphics Primitives}), MeRF (\textit{Memory-Efficient Neural Radiance Field}), Pixel-Nerf, Gnt (\textit{Generalizable Nerf Transformer}) \cite{muller2022, reiser2023merf, yu2021pixelnerfneuralradiancefields, t2023attentionnerfneeds}, entre outros, abordam temas comuns, como: (i) \textit{encoding} espacial de parâmetros; e (ii) uso de arquiteturas híbridas entre Convolutional Neural Networks (CNN) e Vision Transformers (ViT).

Este trabalho tem, portanto, o objetivo de explorar o uso dos diferentes modelos e suas respectivas técnicas, ao aplicá-las em alguns conjuntos de dados com vistas esparsas, visando promover uma análise exploratória entre desempenho computacional e qualidade dos campos gerados. Para avaliar a fidelidade das reconstruções, utilizar-se-ão métricas de qualidade visual amplamente consagradas na literatura, como PSNR (\textit{Peak Signal-to-Noise Ratio}), SSIM (\textit{Structural Similarity Index Measure}) \cite{wang2004image} e LPIPS (\textit{Learned Perceptual Image Patch Similarity}) \cite{zhang2018unreasonable}. Paralelamente, a quantificação do custo operacional basear-se-á em métricas de desempenho sistêmico, incluindo o tempo de convergência no treinamento, a velocidade de inferência (medida em \textit{Frames Per Second} - FPS) e a pegada de memória (\textit{memory footprint}), sobretudo o consumo de VRAM (\textit{Video Random Access Memory}). Em suma, espera-se que esta investigação forneça um panorama consolidado sobre os \textit{trade-offs} entre eficiência computacional e fidelidade visual, estabelecendo diretrizes claras para a viabilização e aplicação prática de modelos de síntese de novas vistas em cenários com restrições severas de dados e de hardware.

% Full papers must respect the page limits defined by the conference.
% Conferences that publish just abstracts ask for \textbf{one}-page texts.

\section{Metodologia} 
\label{sec:metodology}

% The first page must display the paper title, the name and address of the
% authors, the abstract in English and ``resumo'' in Portuguese (``resumos'' are
% required only for papers written in Portuguese). The title must be centered
% over the whole page, in 16 point boldface font and with 12 points of space
% before itself. Author names must be centered in 12 point font, bold, all of
% them disposed in the same line, separated by commas and with 12 points of
% space after the title. Addresses must be centered in 12 point font, also with
% 12 points of space after the authors' names. E-mail addresses should be
% written using font Courier New, 10 point nominal size, with 6 points of space
% before and 6 points of space after.

% The abstract and ``resumo'' (if is the case) must be in 12 point Times font,
% indented 0.8cm on both sides. The word \textbf{Abstract} and \textbf{Resumo},
% should be written in boldface and must precede the text.

\section{Experimentos e Resultados}

% In some conferences, the papers are published on CD-ROM while only the
% abstract is published in the printed Proceedings. In this case, authors are
% invited to prepare two final versions of the paper. One, complete, to be
% published on the CD and the other, containing only the first page, with
% abstract and ``resumo'' (for papers in Portuguese).

\section{Conclusões}

% Section titles must be in boldface, 13pt, flush left. There should be an extra
% 12 pt of space before each title. Section numbering is optional. The first
% paragraph of each section should not be indented, while the first lines of
% subsequent paragraphs should be indented by 1.27 cm.


\bibliographystyle{sbc}
\bibliography{sbc-template}

\end{document}
